\scanpage{37}
% sheet: reused_typed
% hand: typed
% L1 | top centre, in reading orientation
\mspage{-3.9-}

% L2–3 | top
leaves $x$ fixed. Secondly, $3)\handwritten{\Rightarrow}1)$ is an immediate consequence of Proposition 2.1.

% L4–11 | upper and middle
Thirdly, we will prove that $2)\handwritten{\Rightarrow}3)$. By virtue of 2), all the points $x_i$ cannot lie in one hyperplane, hence we can choose $n+1$ of the $x_i$ which are not in one hyperplane, say $x^{(\unsure{0|o})},x^{(1)},\ldots,x^{(n)}$. Without loss of generality, we may assume that $x^{(i)}$ has homogeneous coordinates $x_j^{(i)}=\delta_{ij}$. Now for each $n+k$ between $n+1$ and $m$, let $S_k'$ be the set of integers $i$ such that
\[
D_{0,1,\ldots,\handwritten{\hat}{i},\ldots,n,n+k}\handwritten{\ne}0,
\]
i.e. $x^{(n+k)}$ is not in the hyperplane spanned by $x^{(0)},\ldots,\handwritten{\widehat}{x^{(i)}},\ldots,x^{(n)}$.

% L12–18 | middle
Then I claim that there is no partition of the set $\handwritten{\{}0,1,\ldots,n\handwritten{\}}$ into 2 disjoint subsets $T'$ and $T''$ such that every $S_k'$ is contained in $T'$ or $T''$. For if there were, and if one let $L'$ (resp.$L''$) be the linear subspace defined by $X_i=0$ for all $i\in T'$ (resp. $i\in T''$), then every point $x^{(k)}$ would be in $L'\handwritten{\cup}L''$. It follows immediately that a suitable set of subsets $S_i\handwritten{\subset}S_i'$ is an $R$-partition $R$ and that $x\in\handwritten{\unsure{\mathcal U|U}}_R$. \uline{QED}

% L19–20 | lower
From this Proposition, it is clear that $\handwritten{\mathcal U}_{\mathrm{reg}}$ is precisely the set of pre-stable points of $(\handwritten{\mathbb P}_n)^{m+1}$, in the sense of Chapter 1.

% L21–25 | foot
Our next step is to construct a geometric quotient of $\handwritten{\mathcal U}_{\mathrm{reg}}$ by $\mathrm{PGL}(n)$: let $\handwritten{\mathcal U}_1,\ldots,\handwritten{\mathcal U}_N$ be the open subsets $\handwritten{\mathcal U}_R$ of $\handwritten{\mathcal U}_{\mathrm{reg}}$ and the subsets obtained from these by permuting the coordinates. Let $(Z_i,\phi_i)$ be the geometric quotient of $\handwritten{\mathcal U}_i$ by $\mathrm{PGL}(n)$. For all pairs $i,j$, $\handwritten{\mathcal U}_i\handwritten{\cap}\handwritten{\mathcal U}_j$ is an invariant open subset in $\handwritten{\mathcal U}_i$ and in $\handwritten{\mathcal U}_j$.
% archive sheet number: 36

